\bookchapter{第九章：移動辦公與學習}

\chapterintro{吾昔亦以爲手機辦公不過“湊合”而已，直至某次出差忘攜電腦，強以手機成標書之改與合同之籤。實言之，手機辦公之瓶頸往往不在手機本身，而在君不諳其能之邊界也。}

\booksection{一、文檔之處理：於掌中編輯與創作}

於手機屏上編輯文檔，與電腦上之感確有不同。最大之患乃屏小、指粗、排版易偏。然有數術可使體驗大善：

\begin{itemize}
  \item \textbf{WPS Office}：本土化最善，盡兼容微軟之格式（docx/xlsx/pptx），內置PDF工具甚豐。吾尤愛其"語音備註"之能——閱合同時不便打字，徑對手機言數語，即自動轉爲文字批註附於文中。又WPS內置海量免費模板，制簡歷、彙報PPT時可省甚多之時。
  \item \textbf{Microsoft 365}：若君之公司用正版Office，此乃最佳之選。與桌面版格式兼容幾近完美，無WPS開復雜排版微移之患。深融OneDrive，君於電腦上改至一半之文檔，手機上啓即最新之版。
  \item \textbf{石墨文檔/騰訊文檔}：此二者乃"在線協同"之代表。宜多人同編一文檔，每人所改之處實時而顯。吾嘗以石墨文檔與友共制旅行攻略，二人同填信息於表格，效率遠勝一人獨爲。
\end{itemize}

\textbf{實用之工作流}：吾今之習乃——要緊之件存OneDrive，日常編輯用WPS，團隊協作開石墨文檔。文件既上雲端，再無“忘攜U盤”之尷尬矣。

\booksection{二、筆記之用：自隨手記至知識體系}

記筆記之事，衆人初即欲構繁複之系統，然不三日即棄。吾之議乃先輕後重——先養“記”之習，再慮“管理”之難。

\begin{itemize}
  \item \textbf{隨手記——系統自帶即足}：iOS備忘錄啓之極速，且支持掃描文檔、手寫、錄音。安卓各品牌之便籤亦不遜，鎖屏即可速啓。此等工具宜記臨時之想、購物之單、會議之要。其要在啓之速，思及即可記，無需待App加載。
  \item \textbf{稍後讀——收集信息}：微信之"浮窗"實乃隱藏神器。見佳文未及細閱，點右上角"浮窗"，即縮爲小圓點貼於屏邊，隨時可點開續讀。較轉發至"文件傳輸助手"便利多矣。又，Cubox或Pocket此等稍後讀之器，可跨平臺收集君於各處所見之佳文。
  \item \textbf{構體系——Notion與飛書}：當君覺零散之筆記難覓時，即可上Notion。此物甚靈——君可建"讀書筆記數據庫"，自動以書名、作者、評分組織之；亦可搭"項目看板"，使工作任務之進可視化。飛書之"多維表格"亦類此，且國內訪問更速。吾以Notion理所有寫作素材、學習筆記與項目進度，已用三年餘矣。
\end{itemize}

\textbf{一實用之筆記策}：吾己之習乃“三件套”——臨時之想用系統便籤（秒開），待讀之文用微信浮窗（順手），需長期積累之識放Notion（體系化）。勿初即欲將諸物塞入一器，如此反易棄也。

\booksection{三、掃描與OCR：以紙質之世搬入手機}

手機掃描之事，善用之誠可省卻諸多奔走打印店之勞。吾前次辦簽證需一堆材料之電子版，皆以手機成之。

\begin{itemize}
  \item \textbf{iPhone之用戶}：iOS備忘錄自帶之掃描已甚善。開備忘錄 點相機圖標 選"掃描文稿"，對準文件即自動識邊框而矯透視。拍畢可調色——選"黑白"之式出效最似掃描儀，"灰度"宜有圖之資料。
  \item \textbf{安卓之用戶}：多品牌之相機自帶"文檔掃描"之式。若不得，往應用商店下載Microsoft Lens（免費，且與OneDrive相通）。
  \item \textbf{OCR文字識別}：拍畢之掃描件，若需將其上之文字提取以編輯，即需OCR矣。微信自帶之"提取文字"最簡——長按圖片即可用，識別準確率尚可。若大量處理，可用白描App或全能掃描王。於中文識別，百度之OCR引擎效較佳；WPS之圖片轉文字亦甚準。
  \item \textbf{掃描小術}：置於光線充足之處拍之，最好有臺燈自側面打光，免於文檔上留君之手影。若文件折角，先以手壓平再拍。多頁文件可連掃後合爲一PDF，無需一頁頁獨存。
\end{itemize}

\booksection{四、網盤何以擇：無最好者，惟最適合君者}

擇網盤之事吾思之甚久，終覺實則視君之“主力生態”爲何。

\begin{itemize}
  \item \textbf{蘋果全家桶 iCloud}：同步至使君不覺其存。君於手機上拍之照，數分鐘後電腦上即自有。且蘋果之隱私保護最嚴。然其弊亦顯：跨平臺分享予安卓或Windows用戶較繁，免費僅5GB，擴容需費。
  \item \textbf{Windows主力 OneDrive}：若君用Windows電腦+Office 365，OneDrive乃天作之合。Word中存文檔，自動即存於雲端。1TB空間配Office訂閱甚劃算。手機端App之自動備份照片之能亦佳。
  \item \textbf{國內分享 百度網盤}：實言之，百度網盤之下載限速久爲人詬，然其分享生態誠無人能敵。中文互聯網上大量學習資料、影視資源皆以百度網盤分享。若僅偶下載他人分享之件，不充會員亦可湊合。若君爲重度資源收集者，會員該開即開，所省之時值錢多矣。
  \item \textbf{阿里雲盤/誇克網盤}：此二年乃後起之秀。阿里雲盤上傳下載不限速（目前如是），誇克網盤與誇克瀏覽器深綁，搜資料時即可直存至網盤，甚便。
\end{itemize}

\textbf{吾之組合方案}：iCloud理個人照片與日常文件（隱私爲先），百度網盤存資源與受分享（生態爲先），OneDrive理工作文檔（Office集成爲先）。三網盤各司其職，互不相擾。

\booksection{五、一真實之移動辦公場景}

述吾己之真實經歷。嘗有一次出差至外地，客戶臨時欲改合同條款。吾之工作流如是：

\begin{enumerate}
  \item 於酒店收客戶微信發來之修改意見（文字消息+一張手寫批註之照）。
  \item 啓OneDrive尋合同之原件，以WPS啓而編輯，依客戶微信之消息逐條改條款。
  \item 改畢以WPS之"輸出爲PDF"導出，發至客戶微信羣確認。
  \item 客戶提第二輪之意見，吾以WPS之畫筆之器徑於PDF上圈修改之處，配語音批註以明改動之思。
  \item 確認無誤後，以WPS之"電子簽名"在線籤之，發予對方亦籤。
  \item 終以籤畢之合同PDF存檔於百度網盤與iCloud各一，並於Notion之工作日誌記一筆。
\end{enumerate}

全過程自始至終約四十分鐘，全程未開電腦。實言之，三年前吾斷不信手機能爲此等事，然今誠可。要在預諳此等工具，勿待需時方手忙腳亂下載註冊。

\begin{tcbpractice}{試爲之：搭建君之移動辦公工作流}
\begin{enumerate}
  \item 擇定君之主力文檔App（WPS或Microsoft 365），登錄並關聯雲盤。
  \item 開備忘錄或便籤，以掃描之能拍一份文檔，導出爲PDF發予己。
  \item 擇一網盤App，設爲手機照片自動備份（擇"僅Wi-Fi下備份"以省流量）。
  \item 於Notion或飛書新建一頁，試搭一簡之"工作日誌"模板。
\end{enumerate}

\end{tcbpractice}
